\section{Exact refinement}

Let $(\mathbb G,+)$ be an abelian group and let $F,M_r,T_r,q_{r+1}\in\mathbb G$. At refinement depth $r$, split
\begin{equation}
F=M_r+T_r,
\label{eq:body-tail}
\end{equation}
where $M_r$ is the currently admitted body and $T_r$ the explicit open residue. If a newly validated increment $q_{r+1}$ is transferred from the residue to the body, define
\begin{equation}
M_{r+1}=M_r+q_{r+1},\qquad
T_{r+1}=T_r-q_{r+1}.
\label{eq:refinement}
\end{equation}

\begin{theorem}[T11: exact refinement invariance]
Every transfer in \cref{eq:refinement} preserves the total object:
\begin{equation}
M_{r+1}+T_{r+1}=M_r+T_r=F.
\label{eq:refinement-invariance}
\end{equation}
\end{theorem}

\begin{proof}
Substitution gives
\[
(M_r+q_{r+1})+(T_r-q_{r+1})=M_r+T_r.
\]
\end{proof}

\begin{corollary}[No unexplained correction]
If the admitted body changes by $Q$ while the recorded residue does not change by $-Q$, then the total object changes by exactly $Q$. Thus an unexplained correction is not an invariant refinement.
\end{corollary}

\section{Replay-complete append-only persistence}

\begin{definition}[Replay-complete journal]
A journal $\Wset_k=(w_1,\ldots,w_k)$ is replay-complete relative to the fixed datum $\mathfrak F$ when each entry contains sufficient immutable transition data to determine the next state from the preceding state, and there is a fixed deterministic map
\[
\operatorname{Replay}(\mathfrak F,\Wset_k)=\Ldg_k.
\]
The journal is append-only when $\Wset_k$ is a prefix of $\Wset_{k+1}$ for every $k$.
\end{definition}

\begin{theorem}[T12: persistence by journal prefixes]
If the journal is replay-complete and append-only, then for every $0\le j\le n$ the prefix $\Wset_j$ remains uniquely identifiable inside $\Wset_n$ and
\begin{equation}
\Ldg_j=\operatorname{Replay}(\mathfrak F,\Wset_j).
\label{eq:replay}
\end{equation}
Consequently every earlier state, event, validation route, and status transition recorded in that prefix is reconstructible from the later journal.
\end{theorem}

\begin{proof}
Append-only growth gives the canonical prefix injections
\[
\Wset_0\hookrightarrow\Wset_1\hookrightarrow\cdots\hookrightarrow\Wset_n.
\]
Hence the first $j$ immutable entries of $\Wset_n$ recover $\Wset_j$ exactly. Replay completeness then determines $\Ldg_j$ uniquely from the fixed datum $\mathfrak F$ and that prefix.
\end{proof}

\begin{corollary}[Refutation without erasure]
A later entry may change the active accepted view by recording a typed refutation of an earlier event. Since the earlier journal prefix is not rewritten, both the historical acceptance and the later correction remain reconstructible.
\end{corollary}

\section{Main persistent-record theorem}

\begin{theorem}[T13: invariant-gated persistent record]
Assume:
\begin{enumerate}[label=(\roman*)]
\item the structured initial datum of T1;
\item local invariant admissibility as in T2;
\item external validation as in T3;
\item exact target determination or explicit completion as in T4--T5;
\item any declared path-order bound is checked using T6;
\item formal derivations are admitted only by T7;
\item mandatory closure is controlled by a positive gate margin as in T8;
\item fixed-state readiness is evaluated as in T9;
\item status is assigned by the trichotomy T10;
\item correction obeys T11; and
\item state evolution is recorded in a replay-complete append-only journal as in T12.
\end{enumerate}
Then the resulting sequence $(\Ldg_k)_{k\ge0}$ is a persistent guarded state record satisfying:
\begin{enumerate}[label=(\alph*)]
\item no accepted event is assumed initially;
\item every committed transition obeys its declared invariant bound;
\item no candidate can create validity by self-declaration;
\item no target-changing invisible direction is silently admitted;
\item any declared curvature bound is explicit in the transition data;
\item no derivation is admitted without a licensed acyclic route to its exact target;
\item incomplete mandatory closure cannot commit;
\item bounded gate disturbance is tolerated at the fixed midpoint threshold whenever $\Delta_\Theta>2\eta$;
\item wall-clock relabeling preserving causal order does not alter status;
\item every event is exactly committed, held, or rejected;
\item lawful refinement preserves the total body-plus-residue object; and
\item earlier states remain reconstructible from immutable journal prefixes after later correction.
\end{enumerate}
\end{theorem}

\begin{proof}
Apply T1 to the initial state. T2 controls invariant drift. T3 separates candidate content from validation. T4 and T5 eliminate or expose target-changing blind directions and bound target support. T6 measures path-order residue. T7 licenses finite derivations. T8 supplies exact all-obligation admission and its perturbation margin. T9 identifies intrinsic fixed-state readiness and supplies the predicate used explicitly in T10. T10 gives a total disjoint status rule. T11 preserves the complete object under refinement, and T12 preserves earlier records. The listed properties follow componentwise.
\end{proof}

\begin{remark}[Finite core and extension boundary]
The guarded-record theorem T13 and the finite linear-algebraic target constructions T4--T5 are stated in finite dimension unless explicitly noted.  The Hilbert-space theorem family T14--T17 below extends the target-support, finite-channel, reserve, and projection-defect mechanisms to arbitrary Hilbert spaces under explicit finite-obstruction hypotheses.  Neither layer is an oracle for unrestricted truth.
\end{remark}

\section{Exact finite examples}

The following examples are deliberately small. Their purpose is to expose the theorem wiring without numerical fitting.

\subsection{Target completion}
With $A(x,y)=x$ and $L(x,y)=y$, the vector $(0,1)$ is invisible to $A$ but changes $L$. The quotient $\mathfrak Q(A,L)$ is one-dimensional. The completed observation $(x,y)$ makes the quotient vanish.

\subsection{Curvature}
For the matrices in the T6 example, the path-order residue is exactly $(-1/6,1/6)^{\mathsf T}$. Replacing $B$ by a scalar multiple of $A$ gives zero residue.

\subsection{Gate margin}
For $m=4$ and $S(q)=q_1+q_2+q_3+q_4$, the fully closed state has score $4$ while the largest incomplete score is $3$. Thus $\Delta=1$, any threshold in $(3,4)$ is exact, and every uniform disturbance smaller than $1/2$ preserves a separating threshold.

\subsection{Fixed-state selector}
For
\[
G=\operatorname{diag}(0,i,-2i),
\]
the time average in T9 converges to $\operatorname{diag}(1,0,0)$, projecting exactly onto the zero-action coordinate.

\subsection{Refinement}
If $F=(5,7)$, $M_0=(2,3)$, and $T_0=(3,4)$, then transferring $q_1=(1,2)$ gives $M_1=(3,5)$ and $T_1=(2,2)$; both decompositions sum to $F$ exactly.

\section{Finite-core conclusion}

The central point of the finite core is structural. A persistent record need not be postulated as an external database primitive. It can be derived from a sequence of independent obligations: an initial admissibility structure, an invariant transition law, external validation, exact target determination, explicit path-order curvature, finite proof admission, an all-obligation gate with a positive margin, an intrinsic fixed-state selector, a complete status trichotomy, exact refinement, and replay-complete append-only persistence. The main record theorem is therefore deliberately downstream of T1--T12. The Hilbert-space extension that follows strengthens the target-support layer without changing the finite record theorem.
